\documentclass[11pt]{article}
\usepackage[T1]{fontenc}\usepackage{lmodern,amsmath,amssymb,amsthm}
\usepackage[margin=1in]{geometry}\usepackage[hidelinks]{hyperref}
\newtheorem{theorem}{Theorem}\newtheorem{proposition}[theorem]{Proposition}
\newcommand{\Q}{\mathbb Q}\newcommand{\F}{\mathbb F}\newcommand{\Sha}{\mathrm{Sha}}
\title{Marked elliptic curves, point-isolating covers, and isotypic visibility}
\author{\small \href{mailto:admin@cosmosgeek.org}{admin@cosmosgeek.org}\qquad \href{mailto:b321014@iiit-bh.ac.in}{b321014@iiit-bh.ac.in}}
\date{Working paper for specialist review}
\begin{document}\maketitle
\begin{abstract}
A rational point $P=(x,y)$ on $Y^2=f(X)$ with $y\ne0$ can be isolated,
together with its inverse, on infinitely many genus-two covers whose
Jacobian has the same complex analytic rank as the original curve.
At primes of transposition type, a trace congruence detects whether $P$
is twice a point after reduction. We compare this smooth construction to
pinching pairs of rational points, which removes one complex order per
node. Finally we give a dimension obstruction for killing an order-$p$
torsor inside a geometrically isotypic abelian variety. Under a large
residual-image hypothesis the least possible dimension is exactly $p$.
\end{abstract}
\section{A cover that singles out a pair of points}
Let $f(X)=X^3+aX+b$ be separable and $E:Y^2=f(X)$. Fix
$P=(x,y)\in E(\Q)$ with $y\ne0$. For $d\in\Q^\times$ write
\[
 H_d: v^2=f(x-du^2),\qquad A_d:w^2=d(x-X)f(X).
\]
We take smooth projective models, choosing $(x,0)$ as origin on $A_d$.
\begin{theorem}\label{th:cover}
For infinitely many squareclasses of $d$ one has simultaneously
\[
 H_d(\Q)=\{(0,y),(0,-y)\},\qquad L(A_d,1)\ne0.
\]
For these $d$ the two points map to $\pm P$ on $E$ and
$\operatorname{ord}_{s=1}L(\operatorname{Jac}H_d,s)
=\operatorname{ord}_{s=1}L(E,s)$.
\end{theorem}
\begin{proof}
The sextic has six distinct geometric roots, since $f(x)\ne0$.
The involutions $(u,v)\mapsto(-u,v)$ and $(u,v)\mapsto(-u,-v)$
give degree-two quotients $E$ and $A_d$, via
$(X,Y)=(x-du^2,v)$ and $(X,w)=(x-du^2,duv)$ respectively. The
pullbacks of the two invariant differentials span the independent
$u\,du/v$ and $du/v$ lines. They induce a $\Q$-isogeny
$\operatorname{Jac}H_d\sim E\times A_d$.

The curves $A_d$ are the quadratic twists of the fixed elliptic
curve $A_1$. Infinitely many have nonzero central value
\cite{OS} and hence Mordell--Weil rank zero by the rank-zero
theorem for modular elliptic curves \cite{Wiles}. With finitely
many squareclasses excepted their rational torsion is contained
in $A_d[2]$. Indeed over $\Q(\sqrt d)$ a torsion point maps to a
torsion point on $A_1$ of degree at most two. There are only finitely
many such points by the height comparison and Northcott's theorem;
each point of order greater than two can be anti-invariant for at most
one quadratic character. The two-torsion consists of the four branch
points, namely $(x,0)$ and $(e,0)$ for roots $e$ of $f$. A rational
preimage of $(e,0)$ requires $e\in\Q$ and $d$ in the squareclass of
$x-e$. The two points at infinity on the sextic are rational only
if $-d^3$ is a square. Excluding these finitely many classes leaves
only $u=0$ and proves the assertions by the isogeny factorization.
\end{proof}

\begin{proposition}\label{th:bit}
Outside finitely many odd primes, if $f$ has exactly one root modulo
$p$ then
\[
 \frac{a_p(A_d)-a_p(E)}2\equiv
 \mathbf1_{\overline P\notin 2E(\F_p)}\pmod2 .
\]
\end{proposition}
\begin{proof}
Write $e$ for the rational root and $e',e''$ for its conjugates.
For a cubic with unique rational root $e$, the standard two-isogeny
descent map is $P\mapsto X(P)-e$ modulo squares, with value
$f'(e)$ at $(e,0)$. Its kernel is the image of the dual isogeny.
Over $\F_p$ the isogenous curves have the same order and the dual
kernel has size two, so that image has index two and equals $2E(\F_p)$.
In particular, the unique nonzero rational point of order two is
divisible by two exactly when $f'(e)$ is square, equivalently when
$4\mid\#E(\F_p)$.

Set $c=-df(x)$. In the coordinates $Z=c/(X-x)$ and
$W=cw/(X-x)^2$, the companion has a monic cubic with roots
$z_i=c/(e_i-x)$. Its corresponding derivative at $z_e$ is
\[
 (z_e-z_{e'})(z_e-z_{e''})
 =\frac{c^2 f'(e)}{(e-x)^2(e'-x)(e''-x)}.
\]
Since $(e'-x)(e''-x)=f(x)/(x-e)$ and $f(x)=y^2$, its squareclass
is $f'(e)(x-e)$. Thus divisibility by four of the two group orders
differs exactly when $x-e$ is nonsquare. Both orders are even; taking
their difference and using $a_p=p+1-\#E(\F_p)$ proves the formula.
\end{proof}
The bit is unchanged when either curve is quadratically twisted at a
good unramified prime of this type: changing the sign of an even trace
leaves it unchanged modulo four. In particular one may choose a
rank-zero twist of each curve separately by the same nonvanishing
theorem. The bits measure the original marked class, but do not force
a central derivative to vanish.

\section{The finite-field norm and its exact average}
Suppose $E[2]$ has full $S_3$ Galois image. Fix $d$ independent
classes of rational points modulo twice $E(\Q)$, represented by
$P_i$. Given a good odd prime, let $F$ denote its Frobenius endomorphism,
$D=E(\F_{p^2})$, $A=E(\F_p)$ and $B=\ker(1+F)$ on $D$.
Both $1+F$ and $1-F$ are surjective on geometric points; a solution
to $(1+F)Y=a\in A$ or $(1-F)Y=b\in B$ lies in $D$. Their maps
induce natural isomorphisms
\[
 A/2A \ \xleftarrow{\ 1+F\ }\ D/(A+B)
   \ \xrightarrow{\ 1-F\ }\ B/2B .
\]
For example $(1+F)(A+B)=2A$ and its kernel on $D$ is $B$.
Write $\overline P_i^-$ for the resulting classes in $B/2B$.

\begin{proposition}
The prime mean of $\#((B/2B)/\langle\overline P_1^-,\ldots,
\overline P_d^-\rangle)$ is $1+2^{-d}$. At transposition primes the
bits of Proposition~\ref{th:bit} realize each length-$d$ binary pattern.
\end{proposition}
\begin{proof}
Put $V=E[2]$ and $K=\Q(V)$. Since $V^{C_3}=0$ and
$H^1(C_3,V)=0$, inflation--restriction gives $H^1(S_3,V)=0$.
Thus restriction of the $d$ classes to $G_K$ is independent. The
image of their joint homomorphism is an $S_3$-submodule of $V^d$.
The latter is semisimple and $\operatorname{End}_{S_3}V=\F_2$, so a
proper image would yield a nonzero linear combination restricting
to zero. It follows that the translations in the joint division
extension fill $V^d$. For a fixed linear Frobenius $g$, the reductions
in $V/(g-1)V\simeq E(\F_p)/2E(\F_p)$ are therefore independently
uniform under the finite-group Chebotarev average. If this quotient has
dimension $h$, the expected cardinality after dividing by the $d$
classes is $1+(2^h-1)2^{-d}$: count dual characters vanishing on them.
The mean of $2^h-1$ over $S_3$ is one, by transitivity on
$V\setminus\{0\}$. The event in the second assertion is invariant
under conjugation and has positive density by the same calculation.
See also the general localization method in \cite[\S3]{MR}.
\end{proof}

\section{Pinching rational pairs}
Given $k$ nontorsion rational marks $P_i$ independent modulo torsion,
choose pairwise disjoint $\Q$-rational pairs $(A_i,B_i)$ on $E$
with $B_i-A_i=P_i$; multiples of one point avoid any finite set of
forbidden choices. Identify each pair to make a split node on a
proper curve $C$ with normalization $E$. The normalization sequence
for \'etale cohomology gives
\[
 0\longrightarrow\Q_\ell^k\longrightarrow H^1(C_{\overline\Q},\Q_\ell)
 \longrightarrow H^1(E_{\overline\Q},\Q_\ell)\longrightarrow0.
\]
This also follows at good places by counting points over each finite
residue-field extension. Hence outside a fixed finite bad set
$L^S(H^1(C),s)=\zeta^S(s)^k L^S(E,s)$ and the central order drops by
exactly $k$. The dual homology extension
$0\to V_p E\to W\to\Q_p^k\to0$ is the realization of
$[\mathbb Z^k\to E]$ sending the $i$th node generator to $P_i$,
up to choices of orientation. This is the usual normalization
description of the homological Picard 1-motive \cite{Deligne}.
The connecting map sends the standard basis to their Kummer classes.

For a good prime $p$ the induced finite Selmer group is
\[
 H^1_f(\Q,W)=H^1_f(\Q,V_pE)/\langle\kappa(P_i)\rangle .
\]
Indeed $H^1_f(\Q,\Q_p)=0$ by class field theory: a finite class for
the trivial representation is everywhere unramified. A global lift in
$H^1(\Q,V_pE)$ of a finite class of $W$ is locally finite. Away from
$p$ the extension splits since $E(\Q_l)\otimes\Q_p$ in the rational
Kummer sense vanishes for $l\ne p$; at $p$ the extension is crystalline
and splits after tensoring with $B_{\rm cris}$ as Galois modules.
The latter follows from the crystalline comparison isomorphisms and
exactness of $D_{\rm cris}$ on crystalline representations. These
splittings give the required injectivity on the cohomology tests
defining $H^1_f$ \cite[\S3]{BK}. The kernel on global lifts is exactly
the connecting image. Repeating a mark adds another pole to the
analytic product but leaves this image unchanged.

For one node, its generalized Jacobian is an extension of $E$ by
$\mathbb G_m$ with extension class $\pm P\in E^\vee(\Q)$. A lift
of the modular quotient $h:J_0(N)\to E$ would split the pulled-back
extension, so $h^\vee P$ would be zero. Since $hh^\vee$ is a nonzero
integral multiple of the identity (using the Jacobian polarization),
this is impossible even up to isogeny for nontorsion $P$.

\section{The dimension needed to kill an isotypic torsor}
Call an order-$p$ class $\alpha\in\Sha(E/\Q)$ visible in $B$ with
respect to an embedding $i:E\hookrightarrow B$ if it maps to zero
in $H^1(\Q,B)$.
\begin{theorem}
Assume $E$ is non-CM, $p\ge5$, and the image on $E[p]$ contains
$\mathrm{SL}_2(\F_p)$. If $\alpha\in\Sha(E/\Q)$ has exact order $p$,
the smallest possible $\dim B$ among such visibility embeddings with
$B_{\overline\Q}$ isogenous to a power of $E$ is $p$.
\end{theorem}
\begin{proof}
For an embedding into $B$ of dimension $g<p$ put
$M=\operatorname{Hom}_{\overline\Q}(E,B)$. It is free of rank $g$ by
the isotypic and non-CM hypotheses. Its finite Galois image fixes the
primitive vector $i$; divisibility of $i$ by $n>1$ would place $E[n]$
in the kernel. An element of order $p$ in this image would have a
nontrivial rational constituent of dimension $p-1$ as well as the
fixed line, which is impossible. By Cauchy's theorem the field $K$
defining all the maps has degree prime to $p$. Over $K$ the image on
$E[p]$ still contains $\mathrm{SL}_2(\F_p)$, since this group is
generated by elements of order $p$ and the quotient of images has
order prime to $p$.

Extend $i=f_1$ to a $\mathbb Z$-basis of $M$. The evaluation map
$\varphi:E^g\to B$, $(x_j)\mapsto\sum f_j(x_j)$, is an isogeny.
Its kernel contains no $p$-torsion. Otherwise its submodule of $E[p]^g$
would contain a simple copy of $E[p]$, embedded by a vector of scalars
in $\F_p$ by absolute irreducibility. Integral lifts of those scalars
give a nontrivial modulo-$p$ linear combination of the $f_j$ killing
$E[p]$. Such a map factors through $[p]$, contradicting the choice of
basis in $M$. Thus the isogeny's degree $d$ is prime to $p$. Composing
the inverse isogeny cleared by $d$ with projection to the first
coordinate gives a map $r_K:B\to E$ over $K$ with $r_Ki=[d]$.
Summing its Galois conjugates gives a map over $\Q$ with composite
$[K:\Q]d$, a $p$-adic unit. The map on $p$-primary Weil--Ch\^atelet
classes is consequently injective.

Conversely the period--index equality for locally soluble genus-one
torsors over number fields yields a rational divisor class of degree
$p$ represented by an actual divisor. Riemann--Roch gives a rational
effective representative. Each closed point in its support has degree
divisible by $p$, so the representative is one closed point of degree
$p$. Its residue field $L$ splits $\alpha$. Under the diagonal
$E\to\operatorname{Res}_{L/\Q}E_L$ the induced map on $H^1$ is
restriction by Shapiro's lemma. This is the visibility construction
of Agashe--Stein \cite[Prop.~2.4]{AS}; the target has dimension $p$.
\end{proof}
The lower bound concerns geometrically isotypic targets and the stated
residual representation. General visibility allows other isogeny factors.

\section{Lattices marking a point and a torsor}
There is also an exact minimum for the lattice part of a 1-motive.
Let $C$ be an everywhere locally soluble torsor under $E/\Q$ of prime
period $p$, let $\alpha=[C]$, and fix a nontorsion $P\in E(\Q)$.
Write $G$ for the finite Galois image on $Y$ in notation for invariants. Consider a Galois lattice $Y$ and an equivariant homomorphism
$j:Y\to E(\overline{\Q})$. Require that $P$ be in $j(Y^{G})$
and $\alpha$ be in the image of $H^1(\Q,Y)\to H^1(\Q,E)$.
This means that the complex $[Y\to E]$, with $Y$ in degree $-1$,
quotients by the marked point in degree zero and kills the given
torsor in degree one.

\begin{proposition}
The smallest rank of such a lattice is $p$. If $Y$ must be the
degree-zero divisor lattice on a finite reduced puncture set on $C$
and $j$ the divisor-class map, its smallest rank is $2p-1$.
\end{proposition}
\begin{proof}
Let $G$ be the finite image of Galois on $Y$. Its kernel has no
nonzero continuous homomorphism to the free discrete group $Y$,
so $H^1(\Q,Y)=H^1(G,Y)$ is killed by $|G|$. Hence $p\mid|G|$.
An element of order $p$ acts nontrivially on $Y\otimes\Q$ and contributes
a cyclotomic constituent of dimension $p-1$. The invariant vector
mapping to $P$ gives another dimension.

As in the preceding proof, $C$ has a closed point of degree $p$;
let $D$ be its geometric orbit. Put $I_D=\ker(\mathbb Z[D]\to
\mathbb Z)$ with the degree map, and use the canonical map
$I_D\to\operatorname{Pic}^0(C_{\overline\Q})=E(\overline\Q)$.
By Shapiro's lemma $H^1(\Q,\mathbb Z[D])=0$, so the permutation
sequence identifies $H^1(\Q,I_D)$ with $\mathbb Z/p$. A generator
is represented, for $x\in D$, by $\sigma\mapsto\sigma x-x$.
Its image is the torsor class up to sign. Taking $Y=I_D\oplus\mathbb Z$
and mapping the last generator to $P$ gives rank $p$ and the required
maps. In this construction the induced map
$H^1(\Q,Y)\to H^1(\Q,E)$ is injective, so the corresponding
degree-zero hypercohomology is exactly $E(\Q)/\mathbb Z P$.

For any reduced puncture set, its invariant degree-zero divisor
lattice has rank one less than the number of orbits. Every closed
point of $C$ has degree divisible by $p$ by restriction--corestriction.
Hence the rational mark requires at least two orbits, giving lattice
rank at least $2p-1$. Let $[D]+P$ be the translated degree-$p$ class.
Its Brauer obstruction vanishes locally because $C$ has points over
every completion and therefore vanishes globally by the Brauer
local-global theorem. Riemann--Roch gives an effective rational
representative $D'$, which must itself be a single orbit of size $p$.
It is distinct from $D$ because $P$ has infinite order. The union
$D\cup D'$ gives rank $2p-1$; its invariant difference $D'-D$ maps
to $P$, and the map $I_D\to I_{D\cup D'}$ gives the required
cohomology image by functoriality.
\end{proof}
These two bounds use a specified torsor as input. Banwait, Caro and
Chidambaram \cite{BCC} study categorical minimality for visualizing
abelian varieties, a different minimality requirement.

\begin{thebibliography}{9}\raggedright
\bibitem{OS} K. Ono and C. Skinner, \emph{Non-vanishing of quadratic twists
of modular $L$-functions}, \href{https://uva.theopenscholar.com/files/ken-ono/files/037.pdf}{authors' paper}.
\bibitem{Wiles} A. Wiles, \emph{The Birch and Swinnerton-Dyer conjecture},
\href{https://www.claymath.org/wp-content/uploads/2022/05/birchswin.pdf}{Clay description}.
\bibitem{MR} B. Mazur and K. Rubin, \emph{Ranks of twists of elliptic curves
and Hilbert's tenth problem}, \href{https://arxiv.org/abs/0904.3709}{arXiv:0904.3709}.
\bibitem{Deligne} P. Deligne, \emph{Th\'eorie de Hodge III},
\href{https://doi.org/10.1007/BF02685881}{paper}.
\bibitem{BK} S. Bloch and K. Kato, \emph{$L$-functions and Tamagawa
numbers of motives}, \href{https://math.stanford.edu/~conrad/BSDseminar/refs/BKTamagawa.pdf}{paper}.
\bibitem{AS} A. Agashe and W. Stein, \emph{Visibility of Shafarevich--Tate
groups of abelian varieties}, \href{https://wstein.org/papers/visibility_of_sha/jnt_version.pdf}{paper}.
\bibitem{BCC} B. S. Banwait, J. Caro and S. Chidambaram,
\emph{On the Visibility category of the Shafarevich--Tate group},
\href{https://arxiv.org/abs/2601.21519}{arXiv:2601.21519}.
\end{thebibliography}\end{document}